\section{Step 37 Retain the exponent and reduce both initial values}\label{rsa-s37}

After the supplied frame, copy the exponent into its private object and initialize result one. Reduce the original base into the private base, then reduce the identity into temporary storage and copy it to the result. Check each nested return. The second reduction makes modulus one consistent with exponent zero.

Replace only the body of region \texttt{power\_identity}, retaining its BEGIN and END comments.

\Needspace{36\baselineskip}
\MARSCodeLabel{MIPS assembly}
\begin{minted}[samepage=true]{mars-idea}
    move $s0, $a2
    move $s1, $a3
    lwu $t0, 0($a1)
    sw $t0, 16($sp)
    lwu $t0, 4($a1)
    sw $t0, 20($sp)
    lwu $t0, 8($a1)
    sw $t0, 24($sp)
    lwu $t0, 12($a1)
    sw $t0, 28($sp)
    sd $zero, 32($sp)
    sd $zero, 40($sp)
    li $t0, 1
    sw $t0, 32($sp)
    move $a1, $s0
    move $a2, $sp
    jal rsa_reduce128
    nop
    bnez $v0, pow128_development
    nop
    daddiu $a0, $sp, 32
    move $a1, $s0
    daddiu $a2, $sp, 48
    jal rsa_reduce128
    nop
    bnez $v0, pow128_development
    nop
    ld $t0, 48($sp)
    ld $t1, 56($sp)
    sd $t0, 32($sp)
    sd $t1, 40($sp)
\end{minted}

Use recipe \texttt{s37-1} with the cumulative source through this edit. Stop at \texttt{pow128\_loop} before continuing to the temporary or normal return.

\subsection{Expected observation}
Before the loop, private base 65, exponent 17 and result one occupy distinct objects.

\Needspace{9\baselineskip}
\subsection{Check your understanding}
Why use separate temporary output for these calls?

\begin{itemize}
  \item \textbf{A} The result must share the modulus object.
  \item \textbf{B} Their output must be disjoint from input objects.
  \item \textbf{C} A pointer always needs 32 extra bytes.
\end{itemize}
